% Distilled blueprint for Evans, Chapter 8 (The Calculus of Variations).
% Formal declarations, metadata, and citation identifiers are preserved; exposition and exercises are omitted.

\section{Introduction}
\label{sec:calculus-of-variations-introduction}

\subsection{Basic Ideas}
\label{sec:basic-ideas-calculus-of-variations}

\subsection{First Variation, Euler--Lagrange Equation}
\label{sec:first-variation-euler-lagrange}

Let $U\subset\R^n$ be bounded with smooth boundary, and let
\[
L:\R^n\times\R\times\overline U\to\R
\]
be smooth. For smooth $w$ satisfying $w=g$ on $\partial U$, set
\[
I[w]:=\int_U L(Dw(x),w(x),x)\,dx.
\tag{4}
\]

\begin{theorem}[Euler--Lagrange equation for a minimizer]
\label{thm:euler-lagrange-equation-pde-minimizer}
\dcref{ch8:8.1.2}
\group{Euler--Lagrange Equations}
Any smooth minimizer $u$ of $I[\cdot]$, subject to the boundary condition $u = g$ on $\partial U$, is a solution of the Euler--Lagrange partial differential equation
\[
-\sum_{i=1}^{n}(L_{p_i}(Du, u, x))_{x_i} + L_z(Du, u, x) = 0 \quad \text{in } U.
\]
This is a quasilinear, second-order PDE in divergence form.
\end{theorem}

\begin{proof}
\sketch For $v\in C_c^\infty(U)$, the variations $u+\tau v$ remain admissible for small $\tau$. Differentiating $I[u+\tau v]$ at $\tau=0$ gives the weak Euler--Lagrange identity; integration by parts yields the PDE.
\end{proof}

\begin{example}[Dirichlet's principle]
\label{ex:dirichlets-principle-calculus-of-variations}
\uses{thm:euler-lagrange-equation-pde-minimizer}
\dcref{ch8:8.1-ex-1}
\group{Euler--Lagrange Equations}

Take
\[
L(p, z, x) = \frac{1}{2}|p|^2.
\]
Then $L_{p_i} = p_i$ ($i = 1, \ldots, n$), $L_z = 0$; and so the Euler--Lagrange equation associated with the functional
\[
I[w] := \frac{1}{2}\int_U |Dw|^2 \, dx
\]
is
\[
\Delta u = 0 \quad \text{in } U.
\]
This is \textit{Dirichlet's principle} (see \cref{thm:dirichlets-principle}).
\end{example}

\begin{example}[Generalized Dirichlet's principle]
\label{ex:generalized-dirichlets-principle}
\uses{thm:euler-lagrange-equation-pde-minimizer}
\dcref{ch8:8.1-ex-2}
\group{Euler--Lagrange Equations}

Write
\[
L(p, z, x) = \frac{1}{2}\sum_{i,j=1}^{n} a^{ij}(x) p_i p_j - zf(x),
\]
where $a^{ij} = a^{ji}$ ($i, j = 1, \ldots, n$). Then $L_{p_i} = \sum_{j=1}^{n} a^{ij}(x) p_j$ ($i = 1, \ldots, n$), $L_z = -f(x)$. Hence the Euler--Lagrange equation associated with the functional
\[
I[w] := \int_U \left\{\frac{1}{2}\sum_{i,j=1}^{n} a^{ij} w_{x_i} w_{x_j} - wf\right\} dx
\]
is the divergence structure linear equation
\[
-\sum_{i,j=1}^{n}(a^{ij} u_{x_j})_{x_i} = f \quad \text{in } U.
\]
Uniform ellipticity of the coefficients $a^{ij}$ is the coercivity hypothesis that yields existence of a minimizer.
\end{example}

\begin{example}[Nonlinear Poisson equation]
\label{ex:nonlinear-poisson-equation-calculus-of-variations}
\uses{thm:euler-lagrange-equation-pde-minimizer}
\dcref{ch8:8.1-ex-3}
\group{Euler--Lagrange Equations}

Given a smooth function $f : \R \to \R$, define its antiderivative $F(z) = \int_0^z f(y)\, dy$. Then the Euler--Lagrange equation associated with the functional
\[
I[w] := \int_U \frac{1}{2}|Dw|^2 - F(w)\, dx
\]
is the nonlinear Poisson equation
\[
-\Delta u = f(u) \text{ in } U.
\]
\end{example}

\begin{example}[Minimal surfaces]
\label{ex:minimal-surfaces-calculus-of-variations}
\uses{thm:euler-lagrange-equation-pde-minimizer}
\dcref{ch8:8.1-ex-4}
\group{Euler--Lagrange Equations}

Let
\[
L(p, z, x) = (1 + |p|^2)^{1/2};
\]
so that
\[
I[w] = \int_U (1 + |Dw|^2)^{1/2} \, dx
\]
is the area of the graph of the function $w : U \to \R$. The associated Euler--Lagrange equation is
\[
\sum_{i=1}^{n} \left( \frac{u_{x_i}}{(1 + |Du|^2)^{1/2}} \right)_{x_i} = 0 \text{ in } U.
\tag{10}
\]
This partial differential equation is the \textit{minimal surface equation}. The expression $\Div\left(\frac{Du}{(1+|Du|^2)^{1/2}}\right)$ on the left side of (10) is $n$ times the \textit{mean curvature} of the graph of $u$. Thus a \textit{minimal surface} has zero mean curvature.
\end{example}

\subsection{Second Variation}
\label{sec:second-variation-calculus-of-variations}

\begin{proposition}[Second variation necessary condition]
\label{prop:second-variation-convexity-necessary-condition}
\uses{thm:euler-lagrange-equation-pde-minimizer}
\dcref{ch8:8.1.3}
\group{Euler--Lagrange Equations}

If $u$ is a smooth minimizer of $I[\cdot]$ as in \cref{thm:euler-lagrange-equation-pde-minimizer}, then
\[
\sum_{i,j=1}^n L_{p_i p_j}(Du, u, x)\xi_i \xi_j \geq 0 \quad (\xi \in \R^n, x \in U).
\tag{14}
\]
\end{proposition}

\begin{proof}
\sketch The second variation is nonnegative for every compactly supported $v$. Insert localized zigzag variations whose gradients approach a fixed vector $\xi$ near a chosen point; after shrinking the support, the lower-order terms vanish and yield (14).
\end{proof}

\subsection{Systems}
\label{sec:systems-calculus-of-variations}

\subsubsection{Euler--Lagrange equations}

Let
\[
L:\mathbb{M}^{m\times n}\times\R^m\times\overline U\to\R
\]
be smooth. For smooth $\mathbf w:\overline U\to\R^m$ satisfying
$\mathbf w=\mathbf g$ on $\partial U$, set
\[
I[\mathbf w]:=\int_U L(D\mathbf w(x),\mathbf w(x),x)\,dx.
\tag{15}
\]

\begin{theorem}[Euler--Lagrange system for a minimizer]
\label{thm:euler-lagrange-system-pde-minimizer}
\uses{thm:euler-lagrange-equation-pde-minimizer}
\dcref{ch8:8.1.4-a}
\group{Euler--Lagrange Equations}

Any smooth minimizer $\mathbf{u} = (u^1,\ldots,u^m)$ of $I[\cdot]$ defined by (15), taken among functions equal to $\mathbf{g}$ on $\partial U$, solves the coupled, quasilinear \textit{system} of PDE
\[
-\sum_{i=1}^n \left( L_{p_i^k}(D\mathbf{u}, \mathbf{u}, x) \right)_{x_i} + L_{z^k}(D\mathbf{u}, \mathbf{u}, x) = 0 \quad \text{in } U \quad (k = 1, \ldots, m).
\]
These are the \textit{Euler--Lagrange equations} for the functional $I[\cdot]$ defined by (15).
\end{theorem}

\begin{proof}
\sketch Vary one vector component at a time by compactly supported test functions. Vanishing of the first variation gives each equation of the Euler--Lagrange system.
\end{proof}

\subsubsection{Null Lagrangians}

\begin{definition}[Null Lagrangian]
\label{def:null-lagrangian}
\uses{thm:euler-lagrange-system-pde-minimizer}
\dcref{ch8:8.1.4-def-1}
\group{Null Lagrangians}

The function $L$ is called a \textit{null Lagrangian} if the system of Euler--Lagrange equations
\[
-\sum_{i=1}^{n} (L_{p_i^k}(D\mathbf{u}, \mathbf{u}, x))_{x_i} + L_{z_k}(D\mathbf{u}, \mathbf{u}, x) = 0 \quad \text{in } U \quad (k = 1, \ldots, m)
\tag{17}
\]
is automatically solved by all smooth functions $\mathbf{u} : U \to \R^m$.
\end{definition}

\begin{theorem}[Null Lagrangians and boundary conditions]
\label{thm:null-lagrangians-boundary-conditions}
\uses{def:null-lagrangian}
\dcref{ch8:8.1-thm-1}
\group{Null Lagrangians}

Let $L$ be a null Lagrangian. Assume $\mathbf{u}, \bar{\mathbf{u}}$ are two functions in $C^2(U, \R^m)$ such that
\[
\mathbf{u} = \bar{\mathbf{u}} \quad \text{on } \partial U .
\tag{18}
\]
Then
\[
I[\mathbf{u}] = I[\bar{\mathbf{u}}] .
\tag{19}
\]
\end{theorem}

\begin{proof}
\sketch Join two maps with the same trace by the affine path $u_\tau=u+\tau(v-u)$. The first variation of a null Lagrangian vanishes along the path, so the integral is constant in $\tau$.
\end{proof}

\begin{lemma}[Divergence-free rows of the cofactor matrix]
\label{lem:cofactor-divergence-free}
\dcref{ch8:8.1-lem-1}
\group{Null Lagrangians}
Let $\mathbf{u} : \R^n \to \R^n$ be a smooth function. Then
\[
\sum_{i=1}^n (\cof D\mathbf{u})^k_{i,x_i} = 0 \quad (k = 1, \ldots, n).
\tag{20}
\]
\end{lemma}

\begin{proof}
\sketch Differentiate the cofactor formula. Terms containing mixed second derivatives cancel in pairs by symmetry of derivatives and antisymmetry of the determinant expansion.
\end{proof}

\begin{theorem}[Determinants as null Lagrangians]
\label{thm:determinant-null-lagrangian}
\uses{lem:cofactor-divergence-free, def:null-lagrangian}
\dcref{ch8:8.1-thm-2}
\group{Null Lagrangians}

The determinant function
\[
L(P) = \det P \quad (P \in \mathbb{M}^{n \times n})
\]
is a null Lagrangian.
\end{theorem}

\begin{proof}
\sketch Since $D_P(\det P)=\operatorname{cof}P$, the Euler--Lagrange system for $L(P)=\det P$ is the row-divergence of $\operatorname{cof}D\mathbf u$. The cofactor-divergence identity makes this system vanish identically.
\end{proof}

\subsubsection{Application}

\begin{theorem}[Brouwer's Fixed Point Theorem]
\label{thm:brouwer-fixed-point-theorem}
\uses{thm:null-lagrangians-boundary-conditions, thm:determinant-null-lagrangian}
\dcref{ch8:8.1-thm-3}
\group{Null Lagrangians}

Assume
\[
\mathbf{u} : B(0,1) \to B(0,1)
\]
is continuous, where $B(0,1)$ denotes the closed unit ball in $\R^n$. Then $\mathbf{u}$ has a fixed point; that is, there exists a point $x \in B(0,1)$ with
\[
\mathbf{u}(x) = x.
\]
\end{theorem}

\begin{proof}
\sketch A fixed-point-free map would define a continuous retraction of the ball onto its boundary by radial projection. After smooth approximation, null-Lagrangian invariance gives $\int_B\det D\mathbf w=|B|$, while $|\mathbf w|=1$ forces $D\mathbf w$ to be singular and the same integral to vanish.
\end{proof}

\section{Existence of Minimizers}
\label{sec:existence-of-minimizers}

\subsection{Coercivity, Lower Semicontinuity}
\label{sec:coercivity-lower-semicontinuity}

Fix $1<q<\infty$ and set
\[
I[w]:=\int_U L(Dw(x),w(x),x)\,dx
\tag{1}
\]
for functions with trace $w=g$ on $\partial U$.

\subsubsection{Coercivity}

\begin{definition}[Coercivity condition]
\label{def:coercivity-condition-lagrangian}
\dcref{ch8:8.2.1-def-1}
\group{Existence of Minimizers}
We say the Lagrangian $L$ satisfies the coercivity condition (for the exponent $q$ of (3)) if
\[
\begin{cases}
\text{there exist constants } \alpha > 0, \beta \geq 0 \text{ such that} \\
L(p, z, x) \geq \alpha|p|^q - \beta \\
\text{for all } p \in \R^n, z \in \R, x \in U.
\end{cases}
\tag{4}
\]
\end{definition}

\begin{definition}[Admissible class]
\label{def:admissible-class-sobolev-minimization}
\uses{def:coercivity-condition-lagrangian}
\dcref{ch8:8.2.1-def-2}
\group{Existence of Minimizers}

Set
\[
\mathcal{A} := \{w \in W^{1,q}(U) \mid w = g \text{ on } \partial U \text{ in the trace sense}\}
\]
as the class of \textit{admissible} functions. By (4), $I[w]$ is defined in
$\R\cup\{+\infty\}$ for $w\in\mathcal A$.
\end{definition}

\subsubsection{Lower semicontinuity}

\begin{definition}[Minimizing sequence]
\label{def:minimizing-sequence-calculus-of-variations}
\uses{def:admissible-class-sobolev-minimization}
\dcref{ch8:8.2.1-def-3}
\group{Existence of Minimizers}

We call a sequence $\{u_k\}_{k=1}^{\infty} \subset \mathcal{A}$ satisfying (7), i.e.\ $I[u_k]\to m := \inf_{w\in\mathcal{A}} I[w]$, a \textit{minimizing sequence}.
\end{definition}

\begin{definition}[Weakly lower semicontinuous functional]
\label{def:weakly-lower-semicontinuous-functional}
\dcref{ch8:8.2.1-def-4}
\group{Existence of Minimizers}
We say that a function $I[\cdot]$ is (sequentially) weakly lower semicontinuous on $W^{1,q}(U)$, provided
\[
I[u] \leq \liminf_{k \to \infty} I[u_k]
\]
whenever
\[
u_k \rightharpoonup u \text{ weakly in } W^{1,q}(U).
\]
\end{definition}

\subsection{Convexity}
\label{sec:convexity-lower-semicontinuity-existence}

\begin{theorem}[Weak lower semicontinuity]
\label{thm:weak-lower-semicontinuity-convex-lagrangian}
\uses{prop:second-variation-convexity-necessary-condition, def:weakly-lower-semicontinuous-functional}
\dcref{ch8:8.2-thm-1}
\group{Existence of Minimizers}

Assume that $L$ is bounded below, and in addition the mapping $p \mapsto L(p,z,x)$ is convex, for each $z \in \R, x \in U$. Then $I[\cdot]$ is weakly lower semicontinuous on $W^{1,q}(U)$.
\end{theorem}

\begin{proof}
\sketch Convexity gives a supporting-hyperplane inequality for $L(Du_k,u_k,x)$ relative to a truncated region where $u$ and $Du$ are bounded. Weak convergence handles the linear gradient term, compact Sobolev convergence handles the zeroth-order term, and monotone exhaustion removes the truncation.
\end{proof}

\begin{remark}[On the role of weak convergence in the proof]
\label{rem:weak-lower-semicontinuity-proof-remark}
\uses{thm:weak-lower-semicontinuity-convex-lagrangian}
\dcref{ch8:8.2.2-rem-1}
\group{Existence of Minimizers}

Convexity makes the dependence on $Du_k$ linear after applying the supporting-hyperplane inequality, so weak convergence suffices; weak convergence does not generally imply $Du_k\to Du$ almost everywhere. Strong $L^q$ convergence of $u_k$ removes the need for convexity in $z\mapsto L(p,z,x)$.
\end{remark}

\begin{theorem}[Existence of minimizer]
\label{thm:existence-minimizer-convex-coercive}
\uses{thm:weak-lower-semicontinuity-convex-lagrangian, def:coercivity-condition-lagrangian, def:admissible-class-sobolev-minimization}
\dcref{ch8:8.2-thm-2}
\group{Existence of Minimizers}

Assume that $L$ satisfies the coercivity inequality (4) and is convex in the variable $p$. Suppose also the admissible set $\mathcal{A}$ is nonempty.

Then there exists at least one function $u \in \mathcal{A}$ solving
\[
I[u] = \min_{w \in \mathcal{A}} I[w] .
\]
\end{theorem}

\begin{proof}
\sketch A minimizing sequence is bounded by coercivity, hence has a weakly convergent subsequence in the admissible Sobolev space. The trace condition is weakly closed and weak lower semicontinuity makes the limit a minimizer.
\end{proof}

\begin{theorem}[Uniqueness of minimizer]
\label{thm:uniqueness-minimizer-uniformly-convex}
\uses{thm:existence-minimizer-convex-coercive}
\dcref{ch8:8.2-thm-3}
\group{Existence of Minimizers}

Assume $L=L(p,x)$ is independent of $z$, and that there exists $\theta>0$ such that
\[
\sum_{i,j=1}^n L_{p_i p_j}(p,x)\xi_i\xi_j
\geq \theta|\xi|^2
\quad (p,\xi\in\R^n,\ x\in U).
\tag{30}
\]
Then a minimizer $u \in \mathcal{A}$ of $I[\cdot]$ is unique.
\end{theorem}

\begin{proof}
\sketch If two distinct minimizers existed, uniform convexity applied to their midpoint would lower the energy by a positive multiple of the squared Sobolev distance, a contradiction.
\end{proof}

\subsection{Weak Solutions of the Euler--Lagrange Equation}
\label{sec:euler-lagrange-equation-weak-solutions}

Assume, for some $C>0$,
\[
|L(p,z,x)|\leq C(|p|^q+|z|^q+1),
\tag{35}
\]
and
\[
\begin{cases}
|D_pL(p,z,x)|\leq C(|p|^{q-1}+|z|^{q-1}+1),\\
|D_zL(p,z,x)|\leq C(|p|^{q-1}+|z|^{q-1}+1)
\end{cases}
\tag{36}
\]
for $p\in\R^n$, $z\in\R$, and $x\in U$. The associated boundary-value problem is
\[
\begin{cases}
-\displaystyle\sum_{i=1}^n(L_{p_i}(Du,u,x))_{x_i}+L_z(Du,u,x)=0 & \text{in }U,\\
u=g & \text{on }\partial U.
\end{cases}
\tag{37}
\]

\begin{definition}[Weak solution of the Euler--Lagrange equation]
\label{def:weak-solution-euler-lagrange-scalar}
\dcref{ch8:8.2.3-def-1}
\group{Existence of Minimizers}
We say $u \in \mathcal{A}$ is a \textit{weak solution} of the boundary-value problem (37) for the Euler--Lagrange equation provided
\[
\int_U \sum_{i=1}^n L_{p_i}(Du,u,x) v_{x_i} + L_z(Du,u,x) v \, dx = 0
\]
for all $v \in W_0^{1,q}(U)$.
\end{definition}

\begin{theorem}[Solution of the Euler--Lagrange equation]
\label{thm:weak-solution-euler-lagrange-scalar}
\uses{def:weak-solution-euler-lagrange-scalar}
\dcref{ch8:8.2-thm-4}
\group{Existence of Minimizers}

Assume $L$ verifies the growth conditions (35), (36), and $u \in \mathcal{A}$ satisfies
\[
I[u] = \min_{w \in \mathcal{A}} I[w].
\]
Then $u$ is a weak solution of (37).
\end{theorem}

\begin{proof}
\sketch For $v\in W_0^{1,q}(U)$, differentiate $I[u+\tau v]$ at zero. The growth bounds justify differentiation under the integral, and minimality forces the resulting weak Euler--Lagrange identity to vanish.
\end{proof}

\begin{remark}[Weak solutions of the Euler--Lagrange equation need not be minimizers]
\label{rem:weak-solution-minimizer-joint-convexity}
\uses{thm:weak-solution-euler-lagrange-scalar}
\dcref{ch8:8.2.3-rem-1}
\group{Existence of Minimizers}

Weak solutions of (37) need not minimize $I$. If the joint mapping $(p,z)\mapsto L(p,z,x)$ is convex for each $x$, every weak solution is a minimizer. Indeed, if $u\in\mathcal A$ solves
\[
\begin{cases}
-\sum_{i=1}^n (L_{p_i}(Du, u, x))_{x_i} + L_z(Du, u, x) = 0 & \text{in } U \\
u = g & \text{on } \partial U
\end{cases}
\tag{42}
\]
in the weak sense and select any $w \in \mathcal{A}$. Utilizing the convexity of the mapping $(p, z) \mapsto L(p, z, x)$, we have
\[
L(p, z, x) + D_p L(p, z, x) \cdot (q - p) + D_z L(p, z, x) \cdot (w - z) \leq L(q, w, x).
\]
Let $p = Du(x)$, $q = Dw(x)$, $z = u(x)$, $w = w(x)$ and integrate over $U$:
\[
I[u] + \int_U D_p L(Du, u, x) \cdot (Dw - Du) + D_z L(Du, u, x)(w - u) \, dx \leq I[w].
\]
In view of (42) the second term on the right is zero, and therefore $I[u] \leq I[w]$ for each $w \in \mathcal{A}$.
\end{remark}

\subsection{Systems}
\label{sec:systems-existence-minimizers}

\subsubsection{Convexity}

Set
\[
I[\mathbf w]:=\int_U L(D\mathbf w(x),\mathbf w(x),x)\,dx,
\]
where $L:\mathbb M^{m\times n}\times\R^m\times U\to\R$, and assume
\[
L(P,z,x)\geq\alpha|P|^q-\beta
\quad(P\in\mathbb M^{m\times n},\ z\in\R^m,\ x\in U)
\tag{43}
\]
for constants $\alpha>0$ and $\beta\geq0$.

\begin{definition}[Admissible class for systems]
\label{def:admissible-class-systems}
\dcref{ch8:8.2.4-def-1}
\group{Existence of Minimizers}
Set
\[
\mathcal{A} = \{\mathbf{w} \in W^{1,q}(U; \R^m) \mid \mathbf{w} = \mathbf{g} \text{ on } \partial U \text{ in the trace sense}\},
\]
where $\mathbf{g} : \partial U \to \R^m$ is given.
\end{definition}

\begin{theorem}[Existence of minimizer, systems]
\label{thm:existence-minimizer-systems}
\uses{def:admissible-class-systems, thm:existence-minimizer-convex-coercive}
\dcref{ch8:8.2-thm-5}
\group{Existence of Minimizers}

Assume that $L$ satisfies the coercivity inequality (43) and is convex in the variable $P$. Suppose also the admissible set $\mathcal{A}$ is nonempty.

Then there exists $\mathbf{u} \in \mathcal{A}$ solving
\[
I[\mathbf{u}] = \min_{\mathbf{w} \in \mathcal{A}} I[\mathbf{w}].
\]
\end{theorem}

\begin{proof}
\sketch Coercivity bounds a minimizing sequence in the vector-valued Sobolev space. Extract a weakly convergent subsequence, preserve the trace condition by weak closedness, and apply convex weak lower semicontinuity in the matrix variable.
\end{proof}

\begin{theorem}[Uniqueness of minimizer, systems]
\label{thm:uniqueness-minimizer-systems}
\uses{thm:existence-minimizer-systems}
\dcref{ch8:8.2-thm-6}
\group{Existence of Minimizers}

Assume $L$ does not depend on $z$ and the mapping $P \mapsto L(P, x)$ is uniformly convex. Then a minimizer $\mathbf{u} \in \mathcal{A}$ of $I[\cdot]$ is unique.
\end{theorem}

Assume additionally that, for some $C>0$,
\[
\begin{cases}
|L(P,z,x)|\leq C(|P|^q+|z|^q+1),\\
|D_PL(P,z,x)|\leq C(|P|^{q-1}+|z|^{q-1}+1),\\
|D_zL(P,z,x)|\leq C(|P|^{q-1}+|z|^{q-1}+1)
\end{cases}
\tag{44}
\]
for $P\in\mathbb M^{m\times n}$, $z\in\R^m$, and $x\in U$. The Euler--Lagrange system is
\[
\begin{cases}
-\displaystyle\sum_{i=1}^n(L_{p_k^i}(D\mathbf u,\mathbf u,x))_{x_i}
+L_{z^k}(D\mathbf u,\mathbf u,x)=0 & \text{in }U,\\
u^k=g^k & \text{on }\partial U
\end{cases}
\tag{45}
\]
for $k=1,\ldots,m$.

\begin{definition}[Weak solution of the Euler--Lagrange system]
\label{def:weak-solution-euler-lagrange-system}
\dcref{ch8:8.2.4-def-2}
\group{Existence of Minimizers}
We define $\mathbf{u} \in \mathcal{A}$ to be a \textit{weak solution} of (45) provided
\[
\sum_{k=1}^m \int_U \sum_{i=1}^n L_{p_k^i}(D\mathbf{u}, \mathbf{u}, x)w^k_{x_i} + L_{z^k}(D\mathbf{u}, \mathbf{u}, x)w^k \, dx = 0
\]
for all $\mathbf{w} \in W_0^{1,q}(U; \R^m)$, $\mathbf{w} = (w^1, \ldots, w^m)$.
\end{definition}

\begin{theorem}[Solution of the Euler--Lagrange system]
\label{thm:weak-solution-euler-lagrange-system}
\uses{def:weak-solution-euler-lagrange-system, thm:weak-solution-euler-lagrange-scalar}
\dcref{ch8:8.2-thm-7}
\group{Existence of Minimizers}

Assume $L$ verifies the growth conditions (44) and $\mathbf{u} \in \mathcal{A}$ satisfies
\[
I[\mathbf{u}] = \min_{\mathbf{w} \in \mathcal{A}} I[\mathbf{w}].
\]
Then $\mathbf{u}$ is a weak solution of (45).
\end{theorem}

\begin{proof}
\sketch Differentiate $I[\mathbf u+\tau\mathbf w]$ at zero for arbitrary $\mathbf w\in W_0^{1,q}(U;\R^m)$. The growth bounds justify the derivative, and minimality makes it vanish, giving the weak system.
\end{proof}

\subsubsection{Polyconvexity}

\begin{lemma}[Weak continuity of determinants]
\label{lem:weak-continuity-of-determinants}
\uses{lem:cofactor-divergence-free}
\dcref{ch8:8.2-lem-1}
\group{Existence of Minimizers}

Assume $n < q < \infty$ and
\[
\mathbf{u}_k \rightharpoonup \mathbf{u} \quad \text{weakly in } W^{1,q}(U; \R^n).
\]
Then
\[
\det D\mathbf{u}_k \rightharpoonup \det D\mathbf{u} \quad \text{weakly in } L^{q/n}(U).
\]
\end{lemma}

\begin{proof}
\sketch Proceed by induction on the size of minors. Expand each minor using a cofactor row, integrate by parts with the divergence-free cofactor identity, and combine weak convergence of derivatives with strong local convergence of the maps.
\end{proof}

\begin{definition}[Polyconvex Lagrangian]
\label{def:polyconvex-lagrangian}
\dcref{ch8:8.2.4-def-3}
\group{Existence of Minimizers}
A Lagrangian $L$ is \textit{polyconvex} if
\[
L(P,z,x)=F(P,\det P,z,x)
\tag{51}
\]
for a smooth $F:\mathbb M^{n\times n}\times\R\times\R^n\times\overline U\to\R$ such that, for each fixed $(z,x)$, the joint mapping
\[
(P,r)\longmapsto F(P,r,z,x)
\tag{52}
\]
is convex.
\end{definition}

\begin{theorem}[Lower semicontinuity of polyconvex functionals]
\label{thm:lower-semicontinuity-polyconvex}
\uses{def:polyconvex-lagrangian, lem:weak-continuity-of-determinants, thm:weak-lower-semicontinuity-convex-lagrangian}
\dcref{ch8:8.2-thm-8}
\group{Existence of Minimizers}

Suppose $n < q < \infty$. Assume also $L$ is bounded below and is polyconvex. Then $I[u]$ is weakly lower semicontinuous on $W^{1,q}(U;\R^n)$.
\end{theorem}

\begin{proof}
\sketch Lift $D\mathbf u_k$ to $(D\mathbf u_k,\det D\mathbf u_k)$. Weak continuity of determinants gives weak convergence of the lift, and convex weak lower semicontinuity applies to $F$.
\end{proof}

\begin{theorem}[Existence of minimizers, polyconvex functionals]
\label{thm:existence-minimizer-polyconvex}
\uses{thm:lower-semicontinuity-polyconvex, def:admissible-class-systems}
\dcref{ch8:8.2-thm-9}
\group{Existence of Minimizers}

Assume that $n < q < \infty$, and that $L$ satisfies the coercivity inequality (43) and is polyconvex. Suppose also the admissible set $\mathcal{A}$ is nonempty.

Then there exists $u \in \mathcal{A}$ solving
\[
I[u] = \min_{w \in \mathcal{A}} I[w].
\]
\end{theorem}

\begin{example}[Nonlinear elasticity and polyconvexity]
\label{ex:polyconvexity-nonlinear-elasticity}
\uses{def:polyconvex-lagrangian}
\dcref{ch8:8.2.4-ex-1}
\group{Existence of Minimizers}

For a hyperelastic body in dimension $n=3$, the displacement minimizes
\[
I[\mathbf{w}] := \int_U L(D\mathbf{w}, x)\, dx
\]
over $\mathcal A$. A density of the form
\[
L(P,x)=F(P,\det P,x)
\]
records both deformation gradients and local volume change. See Ball \cite{Ball1977}.
\end{example}

\section{Regularity}
\label{sec:regularity-calculus-of-variations}

For $f\in L^2(U)$, set
\[
I[w]:=\int_U L(Dw)-wf\,dx,
\tag{1}
\]
where $|D_pL(p)|\leq C(|p|+1)$. Its weak Euler--Lagrange equation is
\[
-\sum_{i=1}^n(L_{p_i}(Du))_{x_i}=f\quad\text{in }U,
\tag{3}
\]
equivalently
\[
\int_U\sum_{i=1}^nL_{p_i}(Du)v_{x_i}\,dx=\int_Ufv\,dx
\quad(v\in H_0^1(U)).
\tag{4}
\]

\subsection{Second Derivative Estimates}
\label{sec:second-derivative-estimates-calculus-of-variations}

Assume
\[
|D^2L(p)|\leq C
\tag{5}
\]
and, for some $\theta>0$,
\[
\sum_{i,j=1}^nL_{p_ip_j}(p)\xi_i\xi_j\geq\theta|\xi|^2
\quad(p,\xi\in\R^n).
\tag{6}
\]

\begin{theorem}[Second derivatives for minimizers]
\label{thm:second-derivatives-minimizers}
\uses{def:uniform-ellipticity, thm:interior-h2-regularity, thm:boundary-h2-regularity, thm:difference-quotients-weak-derivatives}
\dcref{ch8:8.3-thm-1}
\group{Regularity}

\begin{enumerate}
\item[(i)] Let $u \in H^1(U)$ be a weak solution of the nonlinear partial differential equation (3), where $L$ satisfies (5), (6). Then
\[
u \in H^2_{\mathrm{loc}}(U).
\]
\item[(ii)] If in addition $u \in H_0^1(U)$ and $\partial U$ is $C^2$, then
\[
u \in H^2(U),
\]
with the estimate
\[
\|u\|_{H^2(U)} \le C\|f\|_{L^2(U)}.
\]
\end{enumerate}
\end{theorem}

\begin{proof}
\sketch Insert difference-quotient variations into the weak Euler--Lagrange equation. Uniform convexity controls the highest difference quotients; bounded derivatives of the Lagrangian and Cauchy inequality control errors. The difference-quotient criterion then yields $H^2$ regularity and the estimate.
\end{proof}

\subsection{Remarks on Higher Regularity}
\label{sec:higher-regularity-remarks}

\section{Constraints}
\label{sec:constraints-calculus-of-variations}

\subsection{Nonlinear Eigenvalue Problems}
\label{sec:nonlinear-eigenvalue-problems}

Let $U$ be bounded and connected with smooth boundary, let $G:\R\to\R$ be smooth, and write $g=G'$. Set
\[
I[w]:=\frac12\int_U|Dw|^2\,dx,
\tag{1}
\]
\[
J[w]:=\int_UG(w)\,dx.
\tag{2}
\]
Assume, for some $C>0$,
\[
|g(z)|\leq C(|z|+1),
\tag{3}
\]
\[
|G(z)|\leq C(|z|^2+1).
\tag{4}
\]

\begin{definition}[Admissible class for a nonlinear eigenvalue problem]
\label{def:admissible-class-nonlinear-eigenvalue}
\dcref{ch8:8.4.1-def-1}
\group{Constrained Minimization}
The admissible class is
\[
\mathcal{A} := \{w \in H_0^1(U) \mid J[w] = 0\}.
\]
\end{definition}

\begin{theorem}[Existence of constrained minimizer]
\label{thm:existence-constrained-minimizer-nonlinear-eigenvalue}
\uses{def:admissible-class-nonlinear-eigenvalue}
\dcref{ch8:8.4-thm-1}
\group{Constrained Minimization}

Assume the admissible set $\mathcal{A}$ is nonempty. Then there exists $u \in \mathcal{A}$ satisfying
\[
I[u] = \min_{w \in \mathcal{A}} I[w].
\]
\end{theorem}

\begin{proof}
\sketch The constraint set is weakly closed because the Sobolev embedding into the constraint space is compact. A bounded minimizing sequence therefore has a weak limit satisfying the constraint, and lower semicontinuity yields a minimizer.
\end{proof}

\begin{theorem}[Lagrange multiplier]
\label{thm:lagrange-multiplier-nonlinear-eigenvalue}
\uses{thm:existence-constrained-minimizer-nonlinear-eigenvalue}
\dcref{ch8:8.4-thm-2}
\group{Constrained Minimization}

Let $u \in \mathcal{A}$ satisfy
\[
I[u] = \min_{w \in \mathcal{A}} I[w].
\tag{9}
\]
Then there exists a real number $\lambda$ such that
\[
\int_U Du \cdot Dv\, dx = \lambda \int_U g(u)v\, dx
\tag{10}
\]
for all $v \in H_0^1(U)$.
\end{theorem}

\begin{proof}
\sketch If the derivative of the constraint is nonzero, correct each variation by a scalar multiple of a fixed transverse direction; constrained minimality makes the energy derivative proportional to the constraint derivative. If it vanishes, then $g(u)=0$ a.e. and the multiplier identity is immediate.
\end{proof}

\begin{remark}[Nonlinear eigenvalue problem interpretation]
\label{rem:nonlinear-eigenvalue-lagrange-multiplier-interpretation}
\uses{thm:lagrange-multiplier-nonlinear-eigenvalue}
\dcref{ch8:8.4.1-rem-1}
\group{Constrained Minimization}

The minimizer $u$ is a weak solution of
\[
\begin{cases}
-\Delta u = \lambda g(u) & \text{in } U \\
u = 0 & \text{on } \partial U,
\end{cases}
\tag{11}
\]
where $\lambda$ is the \textit{Lagrange multiplier} corresponding to the \textit{integral constraint}
\[
J[u] = 0.
\tag{12}
\]
For $u\not\equiv0$, (11) is a \textit{nonlinear eigenvalue problem} for $(u,\lambda)$.
\end{remark}

\subsection{Unilateral Constraints, Variational Inequalities}
\label{sec:unilateral-constraints-variational-inequalities}

For smooth $f$ and obstacle $h:\overline U\to\R$, set
\[
I[w]:=\int_U\left(\frac12|Dw|^2-fw\right)\,dx
\tag{24}
\]
and
\[
\mathcal A:=\{w\in H_0^1(U)\mid w\geq h\text{ a.e. in }U\}.
\tag{25}
\]

\begin{definition}[Admissible class for the obstacle problem]
\label{def:admissible-class-obstacle-problem}
\dcref{ch8:8.4.2-def-1}
\group{Variational Inequalities}
The set $\mathcal A$ in (25) is the admissible class for the obstacle problem with obstacle $h$.
\end{definition}

\begin{theorem}[Existence of minimizer, obstacle problem]
\label{thm:existence-minimizer-obstacle-problem}
\uses{def:admissible-class-obstacle-problem}
\dcref{ch8:8.4-thm-3}
\group{Variational Inequalities}

Assume the admissible set $\mathcal{A}$ is nonempty. Then there exists a unique function $u \in \mathcal{A}$ satisfying
\[
I[u] = \min_{w \in \mathcal{A}} I[w].
\]
\end{theorem}

\begin{proof}
\sketch The obstacle class is convex and weakly closed. Coercivity, weak compactness, and lower semicontinuity give a minimizer; strict convexity of the Dirichlet energy on $H_0^1(U)$ gives uniqueness.
\end{proof}

\begin{theorem}[Variational characterization of minimizer]
\label{thm:variational-inequality-characterization}
\uses{thm:existence-minimizer-obstacle-problem}
\dcref{ch8:8.4-thm-4}
\group{Variational Inequalities}

Let $u \in \mathcal{A}$ be the unique solution of
\[
I[u] = \min_{w \in \mathcal{A}} I[w].
\]
Then
\[
\int_U Du \cdot D(w - u)\, dx \geq \int_U f(w - u)\, dx \quad \text{for all } w \in \mathcal{A}.
\tag{26}
\]
We call (26) a \textit{variational inequality}.
\end{theorem}

\begin{proof}
\sketch For a minimizer, test the segment $u+\tau(v-u)$ and take the one-sided derivative at zero. Conversely, convexity bounds $I[v]-I[u]$ below by this first variation, so the variational inequality implies minimality.
\end{proof}

\begin{remark}[Interpretation of the variational inequality]
\label{rem:obstacle-problem-free-boundary-interpretation}
\uses{thm:variational-inequality-characterization}
\dcref{ch8:8.4.2-rem-1}
\group{Variational Inequalities}

If $\partial U$ is smooth, then $u\in W^{2,\infty}(U)$; see Kinderlehrer--Stampacchia \cite{KinderlehrerStampacchia1980}. Hence
\[
O := \{x \in U \mid u(x) > h(x)\}
\]
is open, and
\[
C := \{x \in U \mid u(x) = h(x)\}
\]
is relatively closed. Variations supported in $O$ give $u\in C^\infty(O)$ and
\[
-\Delta u = f \quad \text{in } O.
\tag{28}
\]
Nonnegative variations in $U$ give
\[
-\Delta u \geq f \quad \text{a.e. in } U.
\tag{29}
\]
Consequently
\[
\begin{cases}
u \geq h,\ -\Delta u \geq f & \text{a.e. in } U,\\
-\Delta u = f & \text{on } U \cap \{u > h\}.
\end{cases}
\tag{30}
\]
The \textit{free boundary} is
\[
F := \partial O \cap U
\]
and does not occur explicitly in (30); see Kinderlehrer--Stampacchia \cite{KinderlehrerStampacchia1980}.
\end{remark}

\subsection{Harmonic Maps}
\label{sec:harmonic-maps-pointwise-constraints}

Set
\[
I[\mathbf w]:=\frac12\int_U|D\mathbf w|^2\,dx
\tag{31}
\]
on
\[
\mathcal A:=\{\mathbf w\in H^1(U;\R^m)\mid
\mathbf w=\mathbf g\text{ on }\partial U,\ |\mathbf w|=1\text{ a.e.}\}.
\tag{32}
\]

\begin{theorem}[Euler--Lagrange equation for harmonic maps]
\label{thm:euler-lagrange-harmonic-maps}
\dcref{ch8:8.4-thm-5}
\group{Constrained Minimization}
Let $\mathbf{u} \in \mathcal{A}$ satisfy
\[
I[\mathbf{u}] = \min_{\mathbf{w} \in \mathcal{A}} I[\mathbf{w}].
\]
Then
\[
\int_U D\mathbf{u} : D\mathbf{v}\, dx = \int_U |D\mathbf{u}|^2 \mathbf{u} \cdot \mathbf{v}\, dx
\tag{33}
\]
for each $\mathbf{v} \in H_0^1(U;\R^m) \cap L^\infty(U;\R^m)$.
\end{theorem}

\begin{proof}
\sketch Use the admissible normalized variations $(\mathbf u+\tau\mathbf v)/|\mathbf u+\tau\mathbf v|$. Their derivative at zero is $\mathbf v-(\mathbf u\cdot\mathbf v)\mathbf u$; substituting it into the first variation and using $D\mathbf u^T\mathbf u=0$ gives (33).
\end{proof}

\begin{remark}[Harmonic map equation and the Lagrange multiplier function]
\label{rem:harmonic-maps-lagrange-multiplier-function}
\uses{thm:euler-lagrange-harmonic-maps}
\dcref{ch8:8.4.3-rem-1}
\group{Constrained Minimization}

Identity (33) says that $\mathbf u=(u^1,\ldots,u^m)$ is a weak solution of
\[
\begin{cases} -\Delta \mathbf{u} = |D\mathbf{u}|^2 \mathbf{u} & \text{in } U \\ \mathbf{u} = \mathbf{g} & \text{on } \partial U. \end{cases}
\tag{34}
\]
The function $\lambda=|D\mathbf u|^2$ is the Lagrange multiplier for the pointwise constraint $|\mathbf u|=1$; unlike an integral constraint, a pointwise constraint produces a multiplier function.
\end{remark}

\subsection{Incompressibility}
\label{sec:incompressibility-constraints}

\subsubsection{Stokes' problem}

Let $U\subset\R^3$ be open, bounded, and simply connected, let $\mathbf f\in L^2(U;\R^3)$, and set
\[
I[\mathbf w]:=\int_U\left(\frac12|D\mathbf w|^2-\mathbf f\cdot\mathbf w\right)\,dx
\]
on
\[
\mathcal A:=\{\mathbf w\in H_0^1(U;\R^3)\mid\Div\mathbf w=0\text{ in }U\}.
\]

\begin{definition}[Admissible class for Stokes' problem]
\label{def:admissible-class-stokes-problem}
\dcref{ch8:8.4.4-def-1}
\group{Incompressible Fluids}
The set $\mathcal A$ is the admissible class for Stokes' problem.
\end{definition}

Let $\mathbf u\in\mathcal A$ be the unique minimizer of $I$ over $\mathcal A$.

\begin{theorem}[Pressure as Lagrange multiplier]
\label{thm:pressure-lagrange-multiplier-stokes}
\uses{def:admissible-class-stokes-problem}
\dcref{ch8:8.4-thm-6}
\group{Incompressible Fluids}

There exists a scalar function $p \in L^2_{\text{loc}}(U)$ such that
\[
\int_U Du : Dv\, dx = \int_U p\, \Div \mathbf{v} + \mathbf{f}\cdot\mathbf{v}\, dx
\tag{39}
\]
for all $\mathbf{v} \in H^1(U;\R^3)$ with compact support within $U$.
\end{theorem}

\begin{proof}
\sketch Minimality gives the weak equation against divergence-free variations. The resulting distribution annihilates all divergence-free test fields, so the de Rham/Helmholtz decomposition represents it as a gradient $Dp$, producing the Stokes system.
\end{proof}

\begin{remark}[Weak formulation of Stokes' problem]
\label{rem:stokes-problem-weak-formulation}
\uses{thm:pressure-lagrange-multiplier-stokes}
\dcref{ch8:8.4.4-rem-1}
\group{Incompressible Fluids}

Identity (39) is the weak formulation of \textit{Stokes' problem}
\[
\begin{cases}
-\Delta \mathbf{u} = \mathbf{f} - Dp & \text{in } U \\
\Div \mathbf{u} = 0 & \text{in } U \\
\mathbf{u} = 0 & \text{on } \partial U.
\end{cases}
\tag{40}
\]
The function $p$ is the \textit{pressure} and arises as a Lagrange multiplier corresponding to the \textit{incompressibility condition} $\Div \mathbf{u} = 0$.
\end{remark}

\subsubsection{Incompressible nonlinear elasticity}

For $L:\mathbb M^{3\times3}\times U\to\R$, set
\[
I[\mathbf w]:=\int_U L(D\mathbf w,x)\,dx.
\]

\begin{definition}[Admissible class, incompressible elasticity]
\label{def:admissible-class-incompressible-elasticity}
\dcref{ch8:8.4.4-def-2}
\group{Nonlinear Elasticity}
\[
\mathcal{A} := \{\mathbf{w} \in W^{1,q}(U; \R^3) \mid \mathbf{w} = \mathbf{g} \text{ on } \partial U,\ \det D\mathbf{w} = 1 \text{ a.e.}\}
\]
for some $q > 3$.
\end{definition}

\begin{theorem}[Minimizers with determinant constraint]
\label{thm:existence-minimizer-determinant-constraint}
\uses{def:admissible-class-incompressible-elasticity, lem:weak-continuity-of-determinants}
\dcref{ch8:8.4-thm-7}
\group{Nonlinear Elasticity}

Assume the mapping
\[
P \mapsto L(P,x)
\]
is convex, and $L$ satisfies the coercivity condition
\[
L(P,x) \geq \alpha|P|^q - \beta \quad (P \in \mathbb{M}^{3 \times 3}, x \in U)
\]
for some $\alpha > 0$, $\beta \geq 0$. Suppose finally $\mathcal{A} \neq \emptyset$.

Then there exists $\mathbf{u} \in \mathcal{A}$ satisfying
\[
I[\mathbf{u}] = \min_{\mathbf{w} \in \mathcal{A}} I[\mathbf{w}].
\]
\end{theorem}

\begin{proof}
\sketch Coercivity bounds a minimizing sequence in the admissible vector Sobolev space. Weak continuity of the determinant preserves the incompressibility constraint, and lower semicontinuity yields a minimizer.
\end{proof}

\begin{remark}[Velocity versus displacement incompressibility conditions]
\label{rem:incompressibility-euler-formula}
\uses{def:admissible-class-stokes-problem, def:admissible-class-incompressible-elasticity}
\dcref{ch8:8.4.4-rem-2}
\group{Nonlinear Elasticity}

The Stokes incompressibility condition is
\[
\Div \mathbf{u} = 0
\tag{52}
\]
whereas nonlinear elasticity uses
\[
\det D\mathbf{u} = 1.
\tag{53}
\]
The first constrains a velocity and the second a displacement. For a velocity field $\mathbf w$, let the particle flow solve
\[
\begin{cases}
\dot{\mathbf{y}}(t) = \mathbf{w}(\mathbf{y}(t), t) & (t \in \R) \\
\mathbf{y}(0) = x.
\end{cases}
\]
Writing $\mathbf y(t)=\mathbf y(t,x)$, the flow is volume preserving when
\[
J(x, t) = \det D_x \mathbf{y}(t, x) = 1 \quad \text{for all } x.
\]
Euler's formula gives
\[
J_t(x, t) = (\Div \mathbf{w})(\mathbf{y}(t, x), t)J(x, t),
\]
the divergence taken with respect to the spatial variables. Hence if $\Div \mathbf{w} = 0$, the flow is volume preserving.
\end{remark}

\section{Critical Points}
\label{sec:critical-points-calculus-of-variations}

\subsection{Mountain Pass Theorem}
\label{sec:mountain-pass-theorem-section}

\subsubsection{Critical points, deformations}

Let $H$ be a real Hilbert space and $I:H\to\R$.

\begin{definition}[Fr\'echet differentiability]
\label{def:frechet-derivative-hilbert-space}
\dcref{ch8:8.5.1-def-1}
\group{Critical Point Theory}
We say $I$ is differentiable at $u \in H$ if there exists $v \in H$ such that
\[
I[w] = I[u] + (v, w - u) + o(\|w - u\|) \quad (w \in H).
\tag{1}
\]
The element $v$, if it exists, is unique. We then write $I'[u] = v$.
\end{definition}

\begin{definition}[The class $C^1(H;\R)$]
\label{def:c1-functional-class-hilbert}
\uses{def:frechet-derivative-hilbert-space}
\dcref{ch8:8.5.1-def-2}
\group{Critical Point Theory}

We say $I$ belongs to $C^1(H;\R)$ if $I'[u]$ exists for each $u \in H$, and the mapping $I' : H \to H$ is continuous.
\end{definition}

\begin{remark}[Simplifying Lipschitz hypothesis]
\label{rem:lipschitz-derivative-simplification}
\uses{def:c1-functional-class-hilbert}
\dcref{ch8:8.5.1-rem-1}
\group{Critical Point Theory}

The simplifying hypothesis used here is
\[
I' : H \to H \text{ is Lipschitz continuous on bounded subsets of } H.
\tag{2}
\]
\end{remark}

Let $\mathcal C$ be the class of functionals $I\in C^1(H;\R)$ satisfying (2), and for $c\in\R$ set
\[
A_c:=\{u\in H\mid I[u]\leq c\},
\qquad
K_c:=\{u\in H\mid I[u]=c,\ I'[u]=0\}.
\]

\begin{definition}[Critical points and critical values]
\label{def:critical-point-critical-value}
\uses{def:frechet-derivative-hilbert-space}
\dcref{ch8:8.5.1-def-3}
\group{Critical Point Theory}

(i) We say $u \in H$ is a \textit{critical point} if $I'[u] = 0$.

(ii) The real number $c$ is a \textit{critical value} if $K_c \neq \emptyset$.
\end{definition}

\begin{definition}[Palais--Smale condition]
\label{def:palais-smale-condition}
\dcref{ch8:8.5.1-def-4}
\group{Critical Point Theory}
A functional $I \in C^1(H;\R)$ satisfies the \textit{Palais--Smale compactness condition} if each sequence $\{u_k\}_{k=1}^\infty \subset H$ such that

(i) $\{I[u_k]\}_{k=1}^\infty$ is bounded

and

(ii) $I'[u_k] \to 0$ in $H$,

is precompact in $H$.
\end{definition}

\begin{theorem}[Deformation Theorem]
\label{thm:deformation-theorem}
\uses{def:palais-smale-condition, def:critical-point-critical-value}
\dcref{ch8:8.5-thm-1}
\group{Critical Point Theory}

Assume $I \in \mathcal{C}$ satisfies the Palais--Smale condition. Suppose also
\[
K_c = \emptyset.
\tag{3}
\]
Then for each sufficiently small $\epsilon > 0$, there exists a constant $0 < \delta < \epsilon$ and a function
\[
\eta \in C([0,1] \times H; H)
\]
such that the mappings
\[
\eta_t(u) = \eta(t, u) \quad (0 \le t \le 1, u \in H)
\]
satisfy
\begin{itemize}
\item[(i)] $\eta_0(u) = u$ \quad $(u \in H)$,
\item[(ii)] $\eta_1(u) = u$ \quad $(u \notin I^{-1}[c - \epsilon, c + \epsilon])$,
\item[(iii)] $I[\eta_t(u)] \le I[u]$ \quad $(u \in H, 0 \le t \le 1)$,
\item[(iv)] $\eta_1(A_{c+\delta}) \subset A_{c-\delta}$.
\end{itemize}
\end{theorem}

\begin{proof}
\sketch Away from critical points in the prescribed energy strip, the Palais--Smale condition gives a uniform lower bound on $\|I^{\prime}\|$. Construct a locally Lipschitz pseudo-gradient vector field and flow points downward; the energy decreases at a uniform rate until the required deformation is achieved.
\end{proof}

\subsubsection{Mountain Pass Theorem}

\begin{theorem}[Mountain Pass Theorem]
\label{thm:mountain-pass-theorem}
\uses{thm:deformation-theorem}
\dcref{ch8:8.5-thm-2}
\group{Critical Point Theory}

Assume $I \in \mathcal{C}$ satisfies the Palais--Smale condition. Suppose also
\begin{itemize}
\item[(i)] $I[0] = 0$,
\item[(ii)] there exist constants $r, a > 0$ such that
\[
I[u] \geq a \text{ if } \|u\| = r,
\]
\item[(iii)] there exists an element $v \in H$ with
\[
\|v\| > r, \quad I[v] \leq 0.
\]
\end{itemize}
Define
\[
\Gamma := \{\mathbf{g} \in C([0,1]; H) \mid \mathbf{g}(0) = 0, \, \mathbf{g}(1) = v\}.
\]
Then
\[
c = \inf_{\mathbf{g} \in \Gamma} \max_{0 \le t \le 1} I[\mathbf{g}(t)]
\]
is a \textit{critical value} of $I$.
\end{theorem}

\begin{proof}
\sketch Assume the mountain-pass level is not critical. The deformation theorem lowers every near-optimal path below that level while fixing its endpoints, contradicting the definition of the minimax value.
\end{proof}

\subsection{Application to Semilinear Elliptic PDE}
\label{sec:semilinear-elliptic-application}

Consider
\[
\begin{cases}
-\Delta u=f(u) & \text{in }U,\\
u=0 & \text{on }\partial U.
\end{cases}
\tag{22}
\]
Assume $f$ is smooth and, for
\[
1<p<\frac{n+2}{n-2},
\]
\[
|f(z)|\leq C(1+|z|^p),
\qquad
|f'(z)|\leq C(1+|z|^{p-1}).
\tag{23}
\]
With $F(z):=\int_0^z f(s)\,ds$, assume for some $\gamma<\frac12$ and $0<a\leq A$ that
\[
0\leq F(z)\leq\gamma f(z)z
\tag{24}
\]
and
\[
a|z|^{p+1}\leq|F(z)|\leq A|z|^{p+1}.
\tag{25}
\]

\begin{remark}[A power-law example]
\label{rem:power-nonlinearity-semilinear-example}
\dcref{ch8:8.5.2-rem-1}
\group{Critical Point Theory}
The power-law PDE
\[
-\Delta u = |u|^{p-1}u
\]
satisfies the hypotheses (23)--(25).
\end{remark}

\begin{theorem}[Existence of a nontrivial solution]
\label{thm:existence-nontrivial-solution-semilinear-elliptic}
\uses{thm:mountain-pass-theorem, thm:lax-milgram-theorem}
\dcref{ch8:8.5-thm-3}
\group{Critical Point Theory}

The boundary-value problem (22) has at least one weak solution $u \not\equiv 0$.
\end{theorem}

\begin{proof}
\sketch The associated energy has a strict local minimum at zero and becomes negative along a sufficiently large multiple of a test function, giving mountain-pass geometry. The growth hypotheses and compact Sobolev embedding verify Palais--Smale; the resulting nonzero critical point satisfies the weak PDE.
\end{proof}

\section{References}
\label{sec:calculus-of-variations-references}
\textbf{Variational methods.} \cite{Giaquinta1983}, \cite{GiaquintaHildebrandt1996}, \cite{Zeidler1985}.
\textbf{Lower semicontinuity and regularity.} \cite{LadyzhenskayaUraltseva1968}, \cite{Giaquinta1983}, \cite{GilbargTrudinger1983}, \cite{GilbargTrudinger1983}.
\textbf{Constraints and inequalities.} \cite{KinderlehrerStampacchia1980}, \cite{DacorognaMoser1990}.
\textbf{Critical points.} \cite{Rabinowitz1986}, \cite{Giaquinta1983}.
